\chapter*{Introduction générale}
\addcontentsline{toc}{chapter}{Introduction générale}
\label{chap:introduction}

% ------------------------------------------------------------------
\section*{Contexte général}

L'économie de la République démocratique du Congo est portée, dans une large mesure, par le
commerce de détail et les activités de négoce exercées par des micro et petites entreprises,
formelles ou informelles. À Kinshasa comme dans les grandes agglomérations provinciales, une
fraction considérable des transactions quotidiennes passe par des boutiques, étals de marché
et points de vente qui n'emploient que quelques personnes et dont les données commerciales —
lorsqu'elles existent — restent consignées sur des supports papier ou dans des classeurs
non structurés. Cette situation limite la capacité des commerçants à analyser leur portefeuille
clients, à identifier les acheteurs à fort potentiel ou à détecter les signes précoces d'un
désengagement commercial.

La littérature sur la numérisation des PME en Afrique subsaharienne souligne que l'adoption
d'outils de gestion de la relation client (CRM) par les très petites entreprises se heurte à
plusieurs obstacles conjugués~: coût des licences, exigences en infrastructure réseau, barrière
linguistique des interfaces et inadaptation des modèles de données aux pratiques locales —
notamment la coexistence du dollar américain (USD) et du franc congolais (CDF) dans les
transactions \citep{Sobantu2025, Yoseph2023}. Ces obstacles ont motivé la conception
d'ELMES-TEKA, un progiciel SaaS multi-tenant développé par ELMES pour répondre aux besoins
spécifiques des entrepreneurs congolais formels et informels.

ELMES-TEKA propose, parmi d'autres fonctionnalités, un module de segmentation de la clientèle
fondé sur l'analyse RFM (Récence, Fréquence, Montant), technique largement documentée dans la
littérature en marketing analytique \citep{Fader2005, Salminen2023}. Ce module, dont la
fonction centrale est \texttt{computeRFM}, classe chaque client dans l'une des cinq catégories
prédéfinies en appliquant une règle de tercilisation sur les trois dimensions : la date du
dernier achat (Récence), le nombre de commandes (Fréquence) et la somme des montants
(Montant). La segmentation ainsi produite est destinée à orienter des décisions commerciales —
identifier les clients à fidéliser, ceux à réactiver, ceux qui risquent de partir — dans un
contexte où les ressources humaines et financières consacrées au marketing sont généralement
limitées.

% ------------------------------------------------------------------
\section*{Problématique}

La mise en production d'un module de segmentation embarqué dans un SaaS commercial soulève
une question de fiabilité qui, dans un contexte de ressources contraintes, ne peut être
esquivée~: dans quelle mesure les sorties du module sont-elles conformes à sa spécification, et
dans quelle mesure résistent-elles aux imprécisions ordinaires des données d'entrée~?

Cette question est d'autant plus pertinente que les données collectées par les commerçants
congolais sont susceptibles de présenter des erreurs de saisie, des doublons ou des valeurs
manquantes dans une proportion non négligeable. Si le module de segmentation produit des
résultats différents selon des artefacts numériques indépendants du comportement réel du
client — tels que des erreurs d'arrondi cumulées ou de légères variations de date —, la
valeur décisionnelle des segments s'en trouve directement affectée. Inversement, si le module
s'avère robuste à ces imprécisions, son déploiement auprès de petits commerçants peut être
justifié avec confiance, malgré la qualité imparfaite des données de saisie.

La littérature sur la validation des méthodes RFM \citep{Rajendran2024, Mingmongkhol2023,
Ray2023} ne propose pas de protocole unifié pour évaluer la conformité algorithmique et la
robustesse aux perturbations d'une implémentation logicielle spécifique. Les études
disponibles se concentrent davantage sur la pertinence prédictive des segments produits que
sur la fidélité de l'exécution à une spécification formelle. Cette lacune justifie une
approche expérimentale par simulation, qui permet de soumettre la fonction à des conditions
contrôlées et reproductibles sans nécessiter l'accès à des données de production réelles.

La question de recherche qui guide ce travail est donc la suivante~: \emph{le module RFM
d'ELMES-TEKA, tel qu'implémenté au commit \texttt{b062a93a7be96d082dee9afe66b9f4f1f9ebecaa},
est-il conforme à sa spécification de référence et robuste aux imprécisions courantes des
données d'entrée, dans les conditions testées par simulation~?}

% ------------------------------------------------------------------
\section*{Objectifs de l'étude}

Ce mémoire poursuit deux objectifs complémentaires, formalisés comme hypothèses de travail.

\paragraph{H1 — Conformité à la référence indépendante.}
L'objectif principal est de mesurer le taux de non-conformité entre les sorties du moteur
TypeScript \texttt{computeRFM} et celles d'une implémentation de référence en Python
calculant les sommes monétaires en centimes entiers. La conformité est évaluée sur
12\,000 essais produits par croisement de tailles de corpus, de profils de distribution et
de scénarios de perturbation. L'hypothèse H1 est formulée ainsi~: le taux de
non-conformité est inférieur à un seuil de tolérance de $2\,\%$, jugé compatible avec
un usage décisionnel dans le contexte de PME informelles congolaises.

\paragraph{H2 — Robustesse aux perturbations.}
L'objectif secondaire est d'évaluer la stabilité des affectations de segment lorsque les
données d'entrée subissent des perturbations contrôlées~: bruit monétaire de $\pm1\,\%$,
modification de dates, introduction de doublons ou suppression de commandes. La robustesse
est mesurée par le taux de changement de segment parmi les clients éloignés des seuils de
tercile, c'est-à-dire les clients dont la classification n'est pas, a priori, à la frontière
d'une transition. L'hypothèse H2 est formulée ainsi~: ce taux de changement, mesuré dans
les conditions du critère H2, a une borne supérieure d'intervalle de confiance Monte-Carlo
à $95\,\%$ strictement inférieure ou égale à $5\,\%$.

% ------------------------------------------------------------------
\section*{Approche épistémologique et méthodologique}

Ce travail s'inscrit dans la tradition de la \emph{Design Science Research} (DSR) telle que
définie par \citet{Hevner2004} et opérationnalisée par \citet{Peffers2008}~: ELMES-TEKA
constitue l'artefact construit pour résoudre un problème réel, et la présente étude se
situe dans l'étape \emph{evaluate} du cycle de vie de cet artefact. La posture adoptée est
pragmatique~: la connaissance produite est jugée à l'aune de son utilité pour les décideurs
et les développeurs de la plateforme, non à l'aune d'une vérité absolue sur le comportement
du marché congolais.

Le raisonnement suivi est hypothético-déductif, conformément à la tradition post-positiviste~:
des hypothèses précises sont formulées avant l'analyse, des critères de décision sont fixés
a priori, et les résultats de la simulation permettent de conclure quant à la compatibilité
ou l'incompatibilité des observations avec ces hypothèses. Ce cadre est cohérent avec
l'approche ADEMP (\emph{Aims}, \emph{Data-generating mechanisms}, \emph{Estimands},
\emph{Methods}, \emph{Performance measures}) proposée par \citet{Morris2019} pour structurer
et rapporter les études de simulation en statistique et en informatique.

La méthode de simulation retenue est Monte-Carlo~: pour chaque cellule du plan factoriel
(taille de corpus × profil de distribution × scénario de perturbation), 100 répétitions
indépendantes sont générées par tirage pseudo-aléatoire à graine fixée, ce qui garantit la
reproductibilité complète des résultats. Les intervalles de confiance à $95\,\%$ sont
construits par loi $t$ à partir des 100 répétitions, puis tronqués à $[0;1]$
\citep{Lakens2013}. L'ensemble des données brutes, des scripts d'analyse et des empreintes
SHA-256 de vérification d'intégrité est consigné dans une archive expérimentale \citep{ViannaFilho2026}.

La comparaison de conformité s'appuie sur une implémentation de référence en Python calculant
les sommes monétaires en centimes entiers, contrairement au moteur TypeScript qui opère en
virgule flottante IEEE\,754 double précision \citep{Haissinsky2012, Ribatet_nd}. Cette
différence arithmétique délibérée permet d'isoler la sensibilité de la tercilisation aux
erreurs d'arrondi cumulées.

% ------------------------------------------------------------------
\section*{Délimitation du travail}

La présente évaluation porte exclusivement sur la fonction \texttt{computeRFM} dans sa
version au commit \texttt{b062a93a}, considérée de façon isolée~: les données de
transactions lui sont soumises directement, sans passer par les couches d'authentification,
la base de données MongoDB, les routes HTTP ou le système multi-tenant d'ELMES-TEKA. Ce
périmètre est une simplification nécessaire pour rendre la campagne de simulation faisable
et reproductible dans le cadre d'un mémoire de master, mais il implique que les résultats
ne constituent pas une évaluation de l'ensemble du SaaS.

Les corpus utilisés dans la simulation sont synthétiques. Ils couvrent quatre tailles
($n \in \{3, 10, 50, 200\}$ clients), trois profils de distribution (\texttt{regular},
\texttt{skewed}, \texttt{ties}) et dix scénarios de perturbation. Ces mécanismes générateurs
n'ont pas été calibrés sur des données de transactions réelles de commerçants congolais. Les
conclusions de conformité et de robustesse sont donc relatives aux distributions synthétiques
testées~; leur transposabilité aux conditions de production réelles reste à établir par une
validation terrain ultérieure.

L'étude ne porte pas sur la pertinence prédictive ou la valeur économique des segments
produits~: la question de savoir si la segmentation RFM en terciles constitue le meilleur
modèle pour prédire les comportements de réachat des clients congolais \citep{Fader2005} est
distincte de la question de conformité algorithmique étudiée ici.

% ------------------------------------------------------------------
\section*{Intérêt du travail}

Ce mémoire apporte une contribution à plusieurs niveaux. Sur le plan pratique, il fournit
aux développeurs d'ELMES-TEKA une mesure quantifiée des performances du module RFM et un
diagnostic précis de la cause des non-conformités observées. Sur le plan méthodologique, il
propose une application du cadre ADEMP à l'évaluation d'une fonction embarquée dans un SaaS
africain, domaine pour lequel les protocoles de validation formelle sont peu documentés.
Sur le plan académique, il illustre comment une posture DSR et une simulation Monte-Carlo
peuvent être mobilisées conjointement pour évaluer un artefact logiciel à ressources
limitées, dans un contexte de données synthétiques contrôlées.

% ------------------------------------------------------------------
\section*{Organisation du mémoire}

Ce mémoire est organisé en quatre chapitres.

\medskip

Le \textbf{chapitre~1} présente les matériels et méthodes~: l'artefact ELMES-TEKA et la
fonction \texttt{computeRFM}, la conception du banc de simulation, le plan factoriel de la
campagne, les neuf cas de test manuels, et le cadre ADEMP appliqué à cette étude.

\medskip

Le \textbf{chapitre~2} est consacré à la revue de la littérature~: les fondements de la
segmentation RFM \citep{Fader2005, Salminen2023}, les enjeux de la gestion commerciale
dans les PME africaines et congolaises \citep{Sobantu2025, Yoseph2023}, les méthodes de
simulation en informatique et en statistique \citep{Morris2019}, et les travaux récents sur
la robustesse et la validation des modèles de segmentation \citep{Rajendran2024,
Mingmongkhol2023, Ray2023}.

\medskip

Le \textbf{chapitre~3} présente les résultats de la campagne de façon factuelle~: le bilan
de conformité sur 12\,000 essais, le premier contre-exemple identifié, les résultats du
critère H2 sur les quatre cellules admissibles, la synthèse des taux de changement par taille
et profil, le comportement particulier du profil \texttt{ties}, et les observations
hors domaine.

\medskip

Le \textbf{chapitre~4} discute les résultats~: l'interprétation du taux de non-conformité de
$1{,}667\,\%$, la cause arithmétique identifiée et les pistes de correction, la signification
d'un intervalle de confiance de largeur nulle pour H2, la délimitation du domaine d'usage
licite de la fonction, et les menaces à la validité interne et externe de l'étude.

\medskip

Une \textbf{conclusion générale} synthétise les réponses aux hypothèses H1 et H2, rappelle
les limites de l'évaluation et propose des perspectives de travaux ultérieurs~: validation
sur données réelles, extension multi-devises, et validation prospective de la pertinence
prédictive des segments.
